% filename: abstract-final.tex
% Base: corrected draft provided by student
% Fixes applied:
%   - "up to 49.6%" → "a 49.6% reduction" (matches Ch7/Ch8 framing)
%   - Added robustness finding (+122% EENS under delayed logistics) — Ch8 §8.2
%   - "can significantly reduce" → factual conditional framing
%   - Closing sentence: removed promotional "demonstrate the potential"; replaced
%     with measured finding + caveat, matching Ch8 conclusion tone
%   - Keywords: removed "safe RL"; added "action masking", "ITU dataset",
%     "expected energy not served", "logistics uncertainty"
%   - "solar–battery–diesel" hyphenation standardised
%   - Minor: "Maskable Proximal Policy Optimisation" spelled out on first use
% ============================================================

\chapter*{Abstract}
\addcontentsline{toc}{chapter}{Abstract}

\begin{tabular}{@{} >{\bfseries}p{0.2\textwidth} p{0.77\textwidth} @{}}
\texorpdfstring{\MakeUppercase{Keywords}}{Keywords} &
telecom tower sites; hybrid energy systems; solar PV; battery storage;
diesel logistics; reinforcement learning; constrained Markov decision process;
inventory management; action masking; expected energy not served
\end{tabular}
\par
\begin{doublespacing}

Reliable and low-carbon energy management for hybrid telecom tower sites
remains a critical challenge in regions with weak or intermittent electricity
supply. Conventional rule-based controllers and offline optimisation tools
such as HOMER Pro rely on fixed thresholds and deterministic assumptions,
making them unable to adapt to stochastic solar generation, grid outages, and
fuel logistics uncertainty.

This thesis studies the operational control of hybrid
solar--battery--diesel telecom tower systems under uncertainty in renewable
generation, grid availability, and diesel delivery. The problem is formulated
as a Constrained Markov Decision Process (CMDP) that explicitly models
battery state-of-charge limits and diesel inventory dynamics, including diesel replenishment ordering and stochastic delivery lead times.

A data-grounded simulation environment is developed using the ITU~2024
Smart Energy Scheduling dataset, comprising operational traces from ten
telecom base station sites at hourly resolution. Reinforcement learning
policies are trained using a Maskable Proximal Policy Optimisation
(\texttt{MaskablePPO}) agent with hard action masking to enforce operational
feasibility constraints at every timestep.

The proposed inventory-aware reinforcement learning policy (RLInv) is
evaluated against rule-based controllers and hierarchical decomposition baselines through
a structured eight-variant ablation study under two logistics scenarios.
Under the trained logistics scenario, RLInv achieves a 49.6\% reduction in
Expected Energy Not Served (EENS) relative to the hierarchical decomposition
baseline, a difference that is statistically significant ($p=0.042$,
Cohen's $d=3.67$). However, under a delayed-delivery scenario, RLInv's EENS
increases by 122\%, the largest degradation among all policies evaluated,
indicating that the learned policy does not generalise robustly beyond its
training logistics distribution.

These results show that joint learning of diesel ordering and energy dispatch
reduces unserved energy relative to decoupled ordering strategies under
matched conditions, while classical inventory heuristics offer greater
structural robustness under logistics stress. Both findings contribute to the
understanding of when inventory-aware reinforcement learning provides
reliable advantage for hybrid telecom energy management.

\end{doublespacing}